\documentclass[10pt,a4paper]{amsart}
\usepackage[margin=19mm]{geometry}
\usepackage{amsmath,amssymb,mathtools}
\usepackage[scr=rsfs,cal=euler]{mathalfa}
\usepackage{xfrac}
\usepackage[unicode,hidelinks,bookmarks=false]{hyperref}
\newcommand{\ie}{i.e.\ }
\newcommand{\cf}{cf.\ }
\newcommand{\resp}{respectively\ }
\newcommand{\Subsection}{\subsection}
\providecommand{\nobreakdash}{\nobreak\hbox{-}}
\setlength{\emergencystretch}{2em}
\multlinegap0pt
\usepackage{bm}
\usepackage{bbm}
\usepackage{dsfont}
\usepackage{xspace}
\usepackage{cancel}
\usepackage{mathtools}
\usepackage{amsthm}
\makeatletter
\@ifundefined{lemma}{\newtheorem{lemma}{Lemma}}{}
\@ifundefined{theorem}{\newtheorem{theorem}{Theorem}}{}
\makeatother
\makeatletter
\newcommand{\RR}{{\mathbb R}}
\expandafter\ifx\csname othm\endcsname\relax
\newenvironment{othm}{
  \em
  \vskip 0.10in
  \refstepcounter{othm}
  \noindent{\bf Theorem\ \theothm.}
}


%%%%%%%%%%%%%%%%% My style %%%%%%%%%%%%%%%%%%%

\newcounter{mysec}
\fi
\newcommand{\wt}[1]{\widetilde{#1}}
\makeatother
\title[Non-spectrality of some piecewise smooth curves and unions o]{Non-spectrality of some piecewise smooth curves and unions of line segments}
\author{Mihail N. Kolountzakis, Chun-Kit Lai}
\date{Source: arXiv:2507.00581v1 (CC BY 4.0)}
\begin{document}
\maketitle

\noindent\emph{Source and licence.} Non-spectrality of some piecewise smooth curves and unions of line segments. Version arXiv:2507.00581v1 (CC BY 4.0), distributed under the Creative Commons Attribution 4.0 International licence, as recorded in the arXiv licence metadata for this exact version and shown on the abstract page https://arxiv.org/abs/2507.00581v1. Full rights notice: licenses/arxiv-2507.00581-CC-BY-4.0.txt.\par\smallskip

\noindent\textit{Editorial scope.} Counted unit: the Theorem whose source label is \texttt{theorem-smooth} in the article (source0.tex lines 894--896, source label \texttt{theorem-smooth}), delivered together with the article's own complete original proof of it (source0.tex lines 898--958, one \texttt{proof} environment of 61 lines). No mathematics is written, repaired or re-derived: statement and proof are the article's own text line for line. Required context retained uncounted: eq-auto-correlation (lines 833--835); eq-Gamma(s,t) (lines 837--839); lemma\_self (lines 847--849); lemma-tranverse (lines 860--862); lemma-ball cover (lines 868--870). Every definition, display and label of those lines is retained unchanged. Omitted from this package and not redistributed: 1--832, 836--836, 840--846, 850--859, 863--867, 871--893, 897--897, 959--1067. Nothing inside a retained excerpt is omitted or altered: every retained body is delivered exactly as the article's lines are written, so no omission receipt of any kind accompanies this package. Citations: the retained excerpts carry no citation command at all, so no bibliography entry of the article is transcribed and no reference is redistributed. Reference adaptation: none was needed and none was made; every reference inside the delivered unit resolves inside it, so no unresolved reference or citation is produced. Printed slips kept literal: no printed word, symbol, hypothesis or number of the counted statement or of its proof was altered. Layout and class: the article's own class is \texttt{amsart} with packages including geometry, amssymb, latexsym, amsfonts, textcomp, fancyhdr, calc, graphicx, pdfpages, amsthm, leftidx, amsmath, amstext, amsxtra; the wrapper is the lane's pinned \texttt{amsart} editorial wrapper at 10pt on A4, which restates the theorem-like environments the retained text needs and adds the article's own macro definitions that the retained text needs (\texttt{\textbackslash RR}, \texttt{\textbackslash wt}), with extra wrapper packages (bm, bbm, dsfont, xspace, cancel, mathtools). The delivered numbering is the unit's own. Assets: no retained span contains a figure environment, a graphics-inclusion command, a PDF-page inclusion, a table or a TikZ picture; no image file is copied into this package, the package declares no asset dependency and no image file is redistributed with it. Licence: version arXiv:2507.00581v1 of this article is distributed under the Creative Commons Attribution 4.0 International licence, as recorded in the arXiv licence metadata for this exact version and shown on the abstract page https://arxiv.org/abs/2507.00581v1. A full rights notice is delivered at licenses/arxiv-2507.00581-CC-BY-4.0.txt. Characters: every retained excerpt is pure ASCII, so the delivered engine, XeLaTeX with its default Latin Modern text font, reports no missing character.
\begin{equation}\label{eq-auto-correlation}
\int f(x)d(\mu\ast\wt{\mu})(x) = \int_0^1\int_0^1 f(\gamma(t)-\gamma(s))~dtds
\end{equation}

\begin{equation}\label{eq-Gamma(s,t)}
\Gamma (s,t) = \gamma(s)-\gamma(t). 
\end{equation}

\begin{lemma}\label{lemma_self}
    Let $\gamma$ be a curve with positive curvature. There exists $\delta_0>0$ such that for all intervals $I$ of length less than $\delta_0$, $\Gamma$ is injective on $(I\times I)\setminus\{(s,s):s\in I\}$.
\end{lemma}

\begin{lemma}\label{lemma-tranverse}
    Let $\gamma$ be a smooth closed curve with positive curvature and let $x_0 = \gamma(s_0) = \gamma(t_0)$ be a transverse intersection point. Then there exist intervals $I_{s_0,t_0} = (s_0-\delta_{s_0,t_0},s_0+\delta_{s_0,t_0})\ni s_0$ and $J_{s_0,t_0}=(t_0-\delta_{s_0,t_0},t_0+\delta_{s_0,t_0})\ni t_0$ such that $\Gamma$ is injective on $I\times J$ and the image $\Gamma(I_{s_0,t_0}\times J_{s_0,t_0})$ covers $B(0, \varepsilon_{s_0,t_0})$ for some $\varepsilon_{s_0,t_0}>0$. 
\end{lemma}

\begin{lemma}\label{lemma-ball cover}
    Let $\gamma$ be a smooth closed curve with positive curvature. Then there exists $\varepsilon_1>0$ such that $\Gamma$ maps  onto $B(0,\varepsilon_1)$. 
\end{lemma}

\begin{theorem}\label{theorem-smooth}
    Let $\gamma: [0,1]\to \RR^2$ be a closed smooth curve of positive curvature and {assume that} $\gamma$ has {finitely many self-intersections, all of them transverse}. Let $\mu$ be the arc-length measure on $\gamma$. Then $\mu\ast\wt\mu$ is absolutely continuous with a smooth density in a {punctured} neighborhood of $0$ in $\RR^2$. 
\end{theorem}

\begin{proof}
We first note that {since the} curve is closed and all derivatives agree on the end-points, we may assume that the domain of the interval is the circle ${\mathbb T}$ where $0$ and $1$ are identified as the same point, so that it is a compact set.


    Let $\delta$ be the minimum of $\delta_0$ in Lemma \ref{lemma_self} and all $\delta_{s_0,t_0}$ in Lemma \ref{lemma-tranverse} so that the conclusion for Lemma \ref{lemma_self} and \ref{lemma-tranverse} holds. Since there are only finitely many transverse intersection points {we have} $\delta>0$.  {The union over} all $t\in {\mathbb T}$ {of} $I_t = (t-\delta/4,t+\delta/4)$ covers ${\mathbb T}$. By compactness, we can find finitely many points $t_1,\cdots t_N$ so that ${\mathbb T}$ is covered by {the union of} $I_j = (t_j-\delta/4,t_j+\delta/4)$, {for} $j=1,\cdots, N$. 

     {We introduce a partition of unity subordinate to this covering.} Namely, we can find smooth functions $\varphi_j{\ge 0}$, for $j = 1,\cdots,N$, such that $\varphi_j$ is supported on $I_j$ and 
     $$
     \sum_{j=1}^N \varphi_j(t) = 1, \forall t\in {\mathbb T}.
     $$
Using (\ref{eq-auto-correlation}) and (\ref{eq-Gamma(s,t)}) and the partition of unity {we obtain}
\begin{equation}\label{eq-partition-unity}
\int f(x)d(\mu\ast\wt{\mu})(x) = \sum_{j=1}^N\sum_{k=1}^N\int_{I_j}\int_{I_k} f(\Gamma(s,t))~\varphi_j(s)\varphi_k(t)~dtds.
\end{equation}
Let $\varepsilon'>0$ be such that if $\gamma(\overline{I_j})\cap \gamma(\overline{I_k}) = \varnothing$, then the distance between the compact sets $\gamma(\overline{I_j}), \gamma(\overline{I_k})$ {satisfies}
$$
\mbox{dist} (\gamma(\overline{I_j}), \gamma(\overline{I_k})) >\varepsilon'>0.
$$
This $\varepsilon'>0$ guarantess that if ${\bf u}\in \Gamma(I_j\times I_k)$ and $\|{\bf u}\|<\varepsilon''$, then $\gamma(\overline{I_j})\cap \gamma(\overline{I_k}) \ne  \varnothing$. 

We now define $\varepsilon$ be the minimum of all $\varepsilon_{s_0,t_0}$ in Lemma {\ref{lemma-tranverse}},  {$\varepsilon_1$} defined in Lemma \ref{lemma-ball cover} and $\varepsilon'$ just defined. 


Let  ${\bf u}\in B(0,\varepsilon)\setminus \{0\}$. Lemma \ref{lemma-ball cover} shows that it is possible to find $s,t$ such that $\Gamma (s,t) = {\bf u}$. We now claim that $\mu\ast\wt{\mu}$ has a smooth density function at $u$. This requires us to show that there exists a smooth function $g$ supported in an open neighborhood of $u$ such that 
$$
\int f(x) d(\mu\ast\wt\mu)(x) = \int f(x) g(x)dx
$$
for all continuous functions $f$ supported in that neighborhood.  From (\ref{eq-partition-unity}) and our definition of ${\varepsilon}$, 
\begin{equation}\label{eq-partition-unity1}
\int f(x)d(\mu\ast\wt{\mu})(x) = \sum_{\{(j,k): {\bf u}\in \Gamma(I_j\times I_k)\}}\int_{I_j}\int_{I_k} f(\Gamma(s,t))~\varphi_j(s)\varphi_k(t)~dtds.
\end{equation}
It suffices to show that each of the integrals above can be written as $
\int f({\bf u}) g_{{j,k}}({\bf u})d{\bf u}$ for some smooth functions $g_{{j,k}}$. There will be three cases:
\begin{enumerate}
    \item[Case 1:] {\bf  Same interval.} $I_j = I_k$.
    \item [Case 2:] {\bf  Adjacent interval.} $I_j \cap I_k\ne \varnothing$.
    \item[Case 3:]  {\bf Far away interval.} $I_j \cap I_k= \varnothing$.
\end{enumerate}

\noindent{\bf Case 1:} 
Consider the function {$F: I_j\times I_j\times B(0,\varepsilon)\to \RR^2$}, defined by 
\begin{equation}\label{eq_F}
F(s,t,{\bf x}) = \Gamma(s,t)- {\bf x}.
\end{equation}
By Lemma \ref{lemma_self}, $\Gamma$ is injective on $I_j\times I_j$ removing the diagonal {for any ${\bf x}$}. Let $(s_0,t_0)$ be such that $\Gamma(s_0,t_0) = {\bf u}$. Then $F(s_0,t_0,{\bf u}) =0$.  Note that the Jacobian of $F$ with respect to $(s,t)$ is equal to $J_{\Gamma}$, which is not zero at $(s_0,t_0).$ By the implicit function theorem, there exists a smooth function $g = (g_1,g_2)$ from a neigborhood of ${\bf u}$, denoted by $U$,  to a neighborhood of $(s_0,t_0)$ such that $s = g_1({\bf x})$, $t = g_2({\bf x})$, $(s_0,t_0) = g({\bf u})$ and $F(g({\bf x}),{\bf x}) = 0$. By the change of variable formula, 
\begin{equation}\label{eq-change-variable}
\int_{I_j}\int_{I_j} f(\Gamma(s,t))~\varphi_j(s)\varphi_j(t)~dtds= \int f({\bf x})\cdot \left(\varphi_j (g_1({\bf x}))\varphi_j (g_2({\bf x})) |J_{\Gamma}(g_1({\bf x}),g_2({\bf x}))|^{-1}\right) ~d{\bf x}
\end{equation}
whenever $f$ is supported on $U$. The function  in the bracket provides the smooth density in a neighborhood of ${\bf u}$. 


\noindent{\bf Case 2:} Suppose that ${\bf u}\in \Gamma(I_j\times I_k)$ and $I_j\cap I_k\ne\varnothing$. Then the length of of the interval $I= I_j\cup I_k$ is at most $\delta<\delta_1$. Lemma \ref{lemma_self} shows that $\Gamma$ is injective on $I_j\times I_k$ removing the diagonal points. We can apply the same argument as in Case 1 on $F: I_j\times I_k\times B(0,\varepsilon)\setminus \{0\} \to \RR^2 $ with the same $F$ in (\ref{eq_F}) to obtain a smooth function. 

\noindent{\bf Case 3:} Suppose that ${\bf u}\in \Gamma(I_j\times I_k)$ and $I_j\cap I_k=\varnothing$. Let us write ${\bf u}  = \Gamma (s_1,t_1)$.   By our choice of $\varepsilon<\varepsilon''$, there must be some $s_0\in I_j$ and $t_0\in I_k$ such that $\gamma(s_0) = \gamma(t_0).$ By our choice of $\delta<\delta_{s_0,t_0}$ $I_j\subset I_{s_0,t_0}$ and $I_k\subset J_{s_0,t_0}$, where $I_{s_0,t_0}$ and $J_{s_0,t_0}$ are defined in Lemma \ref{lemma-tranverse}. Hence, $J_{\Gamma}(s_1,t_1)\ne 0$.  By the implicit function theorem, there exists $g = (g_1,g_2)$ from a neighborhood of ${\bf u}$, denoted by $U$, to a neighborhood of $(s_1,t_1)$ such that 
$$
\Gamma (g({\bf x})) = {\bf x}.
$$
Hence, we have a formula as in (\ref{eq-change-variable}), which provides us a smooth density. 

Combining three cases and (\ref{eq-partition-unity1}), we obtain a smooth density at ${\bf u}\in B(0,\varepsilon)\setminus\{0\}$. Since ${\bf u}$ is arbitrary, we establish a smooth density function for all ${\bf u}\in B(0,\varepsilon)\setminus \{0\}$.
\end{proof}

\end{document}
